% !TeX program = xelatex
% 通用中文数学建模竞赛论文骨架。若赛事提供当年官方模板，请以官方模板为准。
% 电子版默认由本文件生成 PDF，不生成 Word；不得包含承诺书和编号专用页。
\documentclass[UTF8,zihao=-4]{ctexart}

\usepackage[a4paper,top=2.5cm,bottom=2.5cm,left=2.5cm,right=2.5cm]{geometry}
\usepackage{amsmath,amssymb,bm}
\usepackage{booktabs,tabularx,array}
\usepackage{graphicx,float}
\usepackage[font=small,labelsep=quad]{caption}
\usepackage[hidelinks]{hyperref}

\newcolumntype{Y}{>{\raggedright\arraybackslash}X}
\graphicspath{{../results/figures/}{figures/}}
\setlength{\parindent}{2em}
\setlength{\parskip}{0pt}
\setlength{\emergencystretch}{2em}
\allowdisplaybreaks[2]
\raggedbottom

\newcommand{\papertitle}{@@PAPER_TITLE@@}

\begin{document}
\pagestyle{plain}
\pagenumbering{arabic}
\setcounter{page}{1}

\begin{center}
  {\zihao{2}\bfseries \papertitle\par}
\end{center}
\vspace{0.5em}

% 摘要应在正文和全部规范结果核验完成后最后撰写。
\begin{abstract}
待补：按“任务—逐问方法—关键定量结果及条件—验证证据—结论”的顺序撰写，并确保所有数字可追溯。
\end{abstract}

\noindent\textbf{关键词：}待补；待补；待补

% 国赛电子论文以摘要专用页开始；正文从下一页开始。其他赛事按其官方模板适配。
\clearpage

\section{问题重述}

待补：只重述决策目标、输入、约束与输出，不照抄原题。

\section{问题分析}

待补：按题目顺序说明机制、难点、耦合关系、基线、主模型与验证路线。

\section{模型假设、符号与数据处理}

\subsection{事实与模型假设}

待补：区分事实、推断、假设、提案与未知项，另记执行/核验状态；说明假设影响与验证方法。

\subsection{符号说明}

待补：统一给出集合、索引、参数、变量、单位和定义域。

\subsection{数据来源、预处理与质量审计}

待补：说明来源、字段、单位、主键、缺失、异常、变换及信息时序。

% ========== 图、表与符号说明的写法示例（本段整体注释，不影响编译结果）==========
% 用法：把需要的片段复制到正文对应位置，删去行首的 % 后按实际内容修改。
%
% 【插图】宽度写成 \textwidth 的比例，不用 6cm 之类的绝对长度：
% 版心宽度随页边距和模板变化，比例写法能自动适配，避免图片超出版心产生 Overfull \hbox 警告。
% 0.85 是常用取值；若图中文字太小，应重新生成图或拆分，而不是继续调大比例。
% 图片文件放在 results/figures 或 paper/figures 下，由前言的 \graphicspath 解析，不写绝对路径。
% \label 必须紧跟在 \caption 之后，否则交叉引用会指向上一个编号。
%
% \begin{figure}[htbp]
%   \centering
%   \includegraphics[width=0.85\textwidth]{q1_route.pdf}
%   \caption{第 1 问最优路径与基线路径对比（$\lambda=0.3$，随机种子 42）}
%   \label{fig:q1-route}
% \end{figure}
%
% 正文用 \ref 交叉引用，不要手写“如下图”或固定编号：
% 由图~\ref{fig:q1-route} 可见，改进模型在相同约束下将总里程降低 12.4\%。
%
% 【三线表】用 booktabs 的 \toprule/\midrule/\bottomrule，不加竖线，正文表格不用 \hline。
% 列数多或表头较长时优先用 tabularx：表格总宽固定为 \textwidth，模板已定义的 Y 列
% （左对齐、可自动换行的 X 列）分摊剩余宽度，避免长单元格把表格撑出版心。
%
% \begin{table}[htbp]
%   \centering
%   \caption{第 1 问不同模型的求解结果对比}
%   \label{tab:q1-compare}
%   \begin{tabularx}{\textwidth}{l r r Y}
%     \toprule
%     模型 & 目标值 & 运行时间/s & 说明 \\
%     \midrule
%     基线（贪心） & 1854.2 & 0.6 & 可行解，未证明最优 \\
%     主模型（MILP） & 1623.7 & 42.1 & 求解状态 optimal，gap 为 0 \\
%     \bottomrule
%   \end{tabularx}
% \end{table}
%
% 【符号说明表】符号、含义、单位与取值范围一次定义，全文、代码和表图保持一致：
%
% \begin{table}[htbp]
%   \centering
%   \caption{主要符号说明}
%   \label{tab:symbols}
%   \begin{tabularx}{\textwidth}{c l c Y}
%     \toprule
%     符号 & 含义 & 单位 & 取值范围与来源 \\
%     \midrule
%     $x_{ij}$ & 节点 $i$ 到节点 $j$ 的运输量 & t & $x_{ij}\ge 0$，决策变量 \\
%     $c_{ij}$ & 单位运输成本 & 元/t & 附件 1 \\
%     $\lambda$ & 成本与时效的权衡系数 & 无量纲 & $\lambda\in[0,1]$，见灵敏度分析 \\
%     \bottomrule
%   \end{tabularx}
% \end{table}
% ==============================================================================

@@QUESTION_SECTIONS@@

\section{模型检验、敏感性与稳健性}

待补：报告与模型类型匹配的验证、参数扰动、不确定性以及失败案例。

\section{模型评价、局限与结论}

待补：结论强度不得超过证据，说明适用边界、局限和可执行建议。

% 按当年规则填写。2026 年国赛要求本声明位于参考文献之前。
\section*{AI 工具使用声明}

待补：如使用 AI，简述用途并指向支撑材料中的 AI工具使用详情.pdf；如未使用，填写当年规则规定的未使用声明。

\begin{thebibliography}{99}
\bibitem{reference-example}
待补：按赛事规定的格式列出实际引用且已核验的文献。
\end{thebibliography}

\appendix
\section{补充结果与复现说明}

待补：记录运行入口、环境、种子、求解状态、输出路径及必要补充表图。

\section{完整源程序与运行说明}

% 2026 国赛要求附录包含完整、可运行的源程序及必要交互命令；支撑压缩包不能替代此项。
% 按实际使用代码纳入共享模块与各问入口。未用程序时按当年规则明确声明。
待补：按当年规则列出实际使用的完整源程序、共享模块、必要软件交互命令和复现步骤。

\section{支撑材料文件列表}

待补：列出全部可运行源程序、使用的自主数据、必要中间结果，以及适用时的 AI工具使用详情.pdf。

\end{document}
